% =============================================================================
% S4 -- Figures (Figures S1-S12)
%
% The twelve figures of paper/figures/, pulled by \inputfigure so that no
% graphic is duplicated in source. Order is fixed and is what the manuscript
% cites as S1...S12: four diagnostics x three regimes, each diagnostic in the
% order stable, shift, placebo.
%   S1-S3   reliability          S4-S6   coverage by single age
%   S7-S9   joint vs marginal    S10-S12 PIT histograms
% fig-reliability-shift appears twice in the submission: as Figure 1 of the
% manuscript and as Figure S2 here.
% Each float is followed by \clearpage so the printed order is the cited order.
% =============================================================================

\clearpage
\section{Figures}
\label{sec:supp-figures}

This section carries the study's twelve figures: four diagnostics, each shown
for all three regimes in the order stable control, shift (train to 2019, test
2020--2024) and placebo (train to 1913, test 1914--1922). The manuscript prints
one of them, the shift-regime reliability diagram, as its Figure~1; it appears
again here as Figure~\ref{fig:reliability-shift} so that this section is
readable on its own and so that the three regimes of that diagnostic can be
compared side by side. The remaining eleven are cited from the manuscript at
the points where its prose reads them.

All twelve are produced by \pkg{scripts/make\_figures.py} from the same
committed sweep outputs as the tables of Section~S3; the stage refits and
rescores nothing, and the provenance stamp on each figure names its inputs.
Four conventions are shared by every panel and are not repeated in the
individual captions. Every figure is panelled by family, so CBD (M5), which is
fitted on ages 55--99, never averages with a full-age family; its age support
is printed in its panel title. Within a family panel the mechanisms are
compared on the intersection of the (origin, population, sex) cells in which
every mechanism present produced a valid row, and that intersection size is
printed in the panel title. Error rows enter no mean, and the footer of each
figure states how many were excluded. Conformal arms are drawn in a distinct
style throughout---dashed line, open marker, dashed bar edge---because they
construct a single interval at 95\%, so they are read at that level only and
never appear in the PIT figures at all.

\begin{figure}[htbp]\centering
\inputfigure{fig-reliability-stable}
\caption{Reliability diagram, stable control regime: empirical coverage against
nominal level (50, 80 and 95\%) per mechanism, one panel per family, with the
diagonal as perfect calibration. Conformal arms contribute a single point at
95\%, their construction level, in a distinct style. This is the control
counterpart of Figure~1 of the manuscript, and the picture behind the finding
that the coverage split runs through the middle of the neural group and is
already open before any break.}
\label{fig:reliability-stable}
\end{figure}
\clearpage

\begin{figure}[htbp]\centering
\inputfigure{fig-reliability-shift}
\caption{Reliability diagram, shift regime (train to 2019, test 2020--2024,
twenty populations): empirical coverage against nominal level (50, 80 and
95\%) per mechanism, one panel per family, with the diagonal as perfect
calibration. Conformal arms contribute a single point at 95\%, their
construction level, in a distinct style. This figure is also Figure~1 of the
manuscript; it is repeated here beside its stable and placebo counterparts.}
\label{fig:reliability-shift}
\end{figure}
\clearpage

\begin{figure}[htbp]\centering
\inputfigure{fig-reliability-placebo}
\caption{Reliability diagram, placebo regime (train to 1913, test 1914--1922,
eleven populations): empirical coverage against nominal level (50, 80 and
95\%) per mechanism, one panel per family, with the diagonal as perfect
calibration. Conformal arms contribute a single point at 95\%, their
construction level, in a distinct style. Read beside
Figure~\ref{fig:reliability-shift}, it is the picture behind the twin-crises
finding that the pandemic's family split reproduces a century earlier.}
\label{fig:reliability-placebo}
\end{figure}
\clearpage

\begin{figure}[htbp]\centering
\inputfigure{fig-coverage-by-age-stable}
\caption{Empirical 95\% interval coverage by single age, stable control regime:
each line is the mean over cells of the per-age hit rate with horizons pooled,
one panel per family and mechanisms as lines; the horizontal reference is the
nominal 0.95. Solid lines are distributional arms, dashed lines conformal arms
read from their interval bounds. Ages a family does not score are absent---CBD
is fitted on ages 55--99, so its panel is a 55--99 curve and is not comparable
with the full-age families'. This is the ungrouped view of \Hfour, and shows
the old-age concentration of the index-extrapolating families in a period with
no structural break.}
\label{fig:coverage-by-age-stable}
\end{figure}
\clearpage

\begin{figure}[htbp]\centering
\inputfigure{fig-coverage-by-age-shift}
\caption{Empirical 95\% interval coverage by single age, shift regime: each
line is the mean over cells of the per-age hit rate with horizons pooled, one
panel per family and mechanisms as lines; the horizontal reference is the
nominal 0.95. Solid lines are distributional arms, dashed lines conformal arms.
Ages a family does not score are absent (CBD: 55--99). It is the ungrouped
counterpart of the three-band summary in Table~\ref{tab:h4-age}, and shows that
the age profile of the failure is family-specific rather than universal, so
\Hfour\ as registered is not supported for every family.}
\label{fig:coverage-by-age-shift}
\end{figure}
\clearpage

\begin{figure}[htbp]\centering
\inputfigure{fig-coverage-by-age-placebo}
\caption{Empirical 95\% interval coverage by single age, placebo regime: each
line is the mean over cells of the per-age hit rate with horizons pooled, one
panel per family and mechanisms as lines; the horizontal reference is the
nominal 0.95. Solid lines are distributional arms, dashed lines conformal arms.
Ages a family does not score are absent (CBD: 55--99). The 1914--1922 window
loads at younger ages than 2020--2024, so the age profile here is read against
Figure~\ref{fig:coverage-by-age-shift} rather than merged with it.}
\label{fig:coverage-by-age-placebo}
\end{figure}
\clearpage

\begin{figure}[htbp]\centering
\inputfigure{fig-joint-vs-marginal-stable}
\caption{Marginal 95\% coverage against joint path coverage over
$h=1,\dots,H$, stable control regime, as paired bars per mechanism with one
panel per family, on the common cells of each panel (intersection size in the
panel title). Conformal arms are drawn in a distinct style and read from their
interval bounds; the nominal level is marked. The gap between the paired bars
is the quantity \Hthree\ is about---the coverage a mortality paper reports
beside the coverage a liability depends on---and it is open in the control as
well as at the break.}
\label{fig:joint-vs-marginal-stable}
\end{figure}
\clearpage

\begin{figure}[htbp]\centering
\inputfigure{fig-joint-vs-marginal-shift}
\caption{Marginal 95\% coverage against joint path coverage over
$h=1,\dots,H$, shift regime, as paired bars per mechanism with one panel per
family, on the common cells of each panel (intersection size in the panel
title). Conformal arms are drawn in a distinct style and read from their
interval bounds; the nominal level is marked. Joint coverage lies below
marginal in every arm shown, and the copula band---the one mechanism
constructed over the path---is the only mechanism whose joint bar stays near
nominal.}
\label{fig:joint-vs-marginal-shift}
\end{figure}
\clearpage

\begin{figure}[htbp]\centering
\inputfigure{fig-joint-vs-marginal-placebo}
\caption{Marginal 95\% coverage against joint path coverage over
$h=1,\dots,H$, placebo regime, as paired bars per mechanism with one panel per
family, on the common cells of each panel (intersection size in the panel
title). Conformal arms are drawn in a distinct style and read from their
interval bounds; the nominal level is marked. The placebo path spans $H=9$
horizons against the shift regime's $H=5$, so the joint bars here are on a
different scale from Figure~\ref{fig:joint-vs-marginal-shift} and the two are
read within regime.}
\label{fig:joint-vs-marginal-placebo}
\end{figure}
\clearpage

\begin{figure}[htbp]\centering
\inputfigure{fig-pit-hist-stable}
\caption{Ten-bin randomised PIT histograms, stable control regime, averaged
across population--sex cells in a family $\times$ distributional-mechanism
grid; the horizontal line is the uniform reference. Conformal arms are absent
because their PIT is a placeholder. A U shape reads as over-confidence, a hump
as under-confidence, a slope as bias. These are the shapes behind the KS
distances of Table~\ref{tab:pit}.}
\label{fig:pit-hist-stable}
\end{figure}
\clearpage

\begin{figure}[htbp]\centering
\inputfigure{fig-pit-hist-shift}
\caption{Ten-bin randomised PIT histograms, shift regime, averaged across
population--sex cells in a family $\times$ distributional-mechanism grid; the
horizontal line is the uniform reference. Conformal arms are absent because
their PIT is a placeholder. A U shape reads as over-confidence, a hump as
under-confidence, a slope as bias. Read beside
Figure~\ref{fig:pit-hist-stable}, it shows that the departures deepen at the
break without changing shape.}
\label{fig:pit-hist-shift}
\end{figure}
\clearpage

\begin{figure}[htbp]\centering
\inputfigure{fig-pit-hist-placebo}
\caption{Ten-bin randomised PIT histograms, placebo regime, averaged across
population--sex cells in a family $\times$ distributional-mechanism grid; the
horizontal line is the uniform reference. Conformal arms are absent because
their PIT is a placeholder. A U shape reads as over-confidence, a hump as
under-confidence, a slope as bias.}
\label{fig:pit-hist-placebo}
\end{figure}
\clearpage
